\begin{ex}
Hình bình hành có mấy cặp cạnh đối diện song song và bằng nhau?
\shortans{2}
\loigiai{
Có 2 cặp cạnh đối diện song song và bằng nhau.
}
\end{ex}

\begin{ex}
Hình thoi có mấy cặp cạnh đối diện song song?
\shortans{2}
\loigiai{
Có 2 cặp cạnh đối diện song song.
}
\end{ex}

\begin{ex}
Hình thoi có bao nhiêu cạnh bằng nhau?
\shortans{4}
\loigiai{
Hình thoi có 4 cạnh bằng nhau.
}
\end{ex}

\begin{ex}
Hai đường chéo của hình thoi có đặc điểm gì đặc biệt? (Vuông góc với nhau ghi 1, bằng nhau ghi 2).
\shortans{1}
\loigiai{
Vuông góc với nhau (1).
}
\end{ex}

\begin{ex}
Hai đường chéo của hình thoi cắt nhau tại điểm nào của mỗi đường? (Trung điểm ghi 1, đầu mút ghi 2).
\shortans{1}
\loigiai{
Tại trung điểm của mỗi đường (1).
}
\end{ex}

\begin{ex}
Hình vuông có phải là một hình thoi đặc biệt không? (Phải ghi 1, không ghi 0).
\shortans{1}
\loigiai{
Hình vuông là hình thoi đặc biệt có 4 góc vuông.
}
\end{ex}

\begin{ex}
Hình vuông có phải là một hình bình hành đặc biệt không? (Phải ghi 1, không ghi 0).
\shortans{1}
\loigiai{
Hình vuông là hình bình hành đặc biệt.
}
\end{ex}

\begin{ex}
Một hình thoi có cạnh dài 8 cm. Chu vi của hình thoi đó là bao nhiêu xăng-ti-mét?
\shortans{32}
\loigiai{
$8 \times 4 = 32$ cm.
}
\end{ex}

\begin{ex}
Một hình thoi có chu vi 48 cm. Cạnh của hình thoi đó dài bao nhiêu xăng-ti-mét?
\shortans{12}
\loigiai{
$48 : 4 = 12$ cm.
}
\end{ex}

\begin{ex}
Một hình bình hành có độ dài hai cạnh liên tiếp là 12 cm và 8 cm. Chu vi hình bình hành đó là bao nhiêu xăng-ti-mét?
\shortans{40}
\loigiai{
$(12 + 8) \times 2 = 40$ cm.
}
\end{ex}

\begin{ex}
Một hình bình hành có chu vi 50 cm, độ dài một cạnh là 15 cm. Độ dài cạnh kia là bao nhiêu xăng-ti-mét?
\shortans{10}
\loigiai{
Nửa chu vi = 25 cm. Cạnh kia = $25 - 15 = 10$ cm.
}
\end{ex}

\begin{ex}
Một hình bình hành có độ dài đáy là 15 cm, chiều cao tương ứng là 8 cm. Diện tích của hình bình hành đó là bao nhiêu xăng-ti-mét vuông?
\shortans{120}
\loigiai{
$15 \times 8 = 120\text{ cm}^2$.
}
\end{ex}

\begin{ex}
Một hình bình hành có diện tích $96\text{ cm}^2$, độ dài đáy là 12 cm. Chiều cao của hình bình hành đó là bao nhiêu xăng-ti-mét?
\shortans{8}
\loigiai{
$96 : 12 = 8$ cm.
}
\end{ex}

\begin{ex}
Một hình thoi có độ dài hai đường chéo là 10 cm và 6 cm. Diện tích của hình thoi đó là bao nhiêu xăng-ti-mét vuông?
\shortans{30}
\loigiai{
$10 \times 6 : 2 = 30\text{ cm}^2$.
}
\end{ex}

\begin{ex}
Một hình thoi có độ dài hai đường chéo là 14 cm và 8 cm. Diện tích của hình thoi đó là bao nhiêu xăng-ti-mét vuông?
\shortans{56}
\loigiai{
$14 \times 8 : 2 = 56\text{ cm}^2$.
}
\end{ex}

\begin{ex}
Một hình thoi có diện tích $40\text{ cm}^2$, độ dài một đường chéo là 8 cm. Độ dài đường chéo kia là bao nhiêu xăng-ti-mét?
\shortans{10}
\loigiai{
$40 \times 2 : 8 = 10$ cm.
}
\end{ex}

\begin{ex}
Một mảnh đất hình bình hành có đáy 24 m, chiều cao 10 m. Diện tích mảnh đất là bao nhiêu mét vuông?
\shortans{240}
\loigiai{
$24 \times 10 = 240\text{ m}^2$.
}
\end{ex}

\begin{ex}
Một chiếc khăn quàng hình thoi có hai đường chéo dài 30 cm và 20 cm. Diện tích chiếc khăn là bao nhiêu xăng-ti-mét vuông?
\shortans{300}
\loigiai{
$30 \times 20 : 2 = 300\text{ cm}^2$.
}
\end{ex}

\begin{ex}
Hình thoi có mấy góc nhọn và mấy góc tù? (Tổng số góc nhọn và góc tù là bao nhiêu?).
\shortans{4}
\loigiai{
Gồm 2 góc nhọn và 2 góc tù, tổng có 4 góc.
}
\end{ex}

\begin{ex}
Hình bình hành có bao nhiêu góc?
\shortans{4}
\loigiai{
Có 4 góc.
}
\end{ex}

\begin{ex}
Một hình thoi có chu vi 60 cm. Độ dài mỗi cạnh của hình thoi là bao nhiêu xăng-ti-mét?
\shortans{15}
\loigiai{
$60 : 4 = 15$ cm.
}
\end{ex}

\begin{ex}
Một hình thoi có hai đường chéo bằng nhau và bằng 10 cm. Diện tích hình thoi đó là bao nhiêu xăng-ti-mét vuông?
\shortans{50}
\loigiai{
$10 \times 10 : 2 = 50\text{ cm}^2$.
}
\end{ex}

\begin{ex}
Một hình bình hành có diện tích $150\text{ m}^2$, chiều cao là 10 m. Độ dài đáy tương ứng là bao nhiêu mét?
\shortans{15}
\loigiai{
$150 : 10 = 15$ m.
}
\end{ex}

\begin{ex}
Hình nào có 4 cạnh bằng nhau và 2 đường chéo vuông góc với nhau? (Hình thoi ghi 1, hình chữ nhật ghi 2).
\shortans{1}
\loigiai{
Hình thoi (1).
}
\end{ex}

\begin{ex}
Hai góc đối diện của hình bình hành có bằng nhau không? (Có ghi 1, không ghi 0).
\shortans{1}
\loigiai{
Hai góc đối diện bằng nhau (1).
}
\end{ex}
